\section{Fields and Vector Spaces}

\subsection{Q1. Unique additive identity}

\begin{promptbox}
State and prove that $\zero\in V$ is the unique additive identity.
\end{promptbox}

\begin{coursebox}
\textbf{Statement.}
There is a unique vector $\zero\in V$ such that $\vv+\zero=\vv$ for every $\vv\in V$.

\textbf{Existence} is an axiom of a vector space.

\textbf{Uniqueness.}
Suppose $\vec{e}\in V$ also satisfies $\vv+\vec{e}=\vv$ for every $\vv\in V$.
Taking $\vv=\zero$ and using that $\zero$ is an additive identity gives
\[
\vec{e}
=\zero+\vec{e}
=\zero.
\]
Thus the additive identity is unique.
\end{coursebox}

\begin{laterbox}
The same one-line uniqueness argument works in any abelian group, so it is not
special to vector spaces. After you know that $V\cong\F^n$ when $\dim V=n<\infty$,
uniqueness of $\zero$ in $V$ is equivalent to uniqueness of the zero tuple in
$\F^n$, but that is a strictly heavier route and hides that the claim is purely
algebraic.
\end{laterbox}

\subsection{Q2. $0\vv=\zero$}

\begin{promptbox}
Prove that $0\vv=\zero$.
\end{promptbox}

\begin{coursebox}
Let $\vv\in V$. Then
\[
0\vv=(0+0)\vv=0\vv+0\vv
\]
by distributivity. Add the additive inverse $-(0\vv)$ to both sides:
\[
\zero=0\vv.
\]
(Equivalently: $0\vv+0\vv=0\vv$ and Q4 below, or cancellation.)
\end{coursebox}

\begin{laterbox}
In coordinates relative to any basis $\cB$, we have $[0\vv]_{\cB}=0[\vv]_{\cB}=\zero$
in $\F^n$, hence $0\vv=\zero$. This uses the coordinate isomorphism, which is
later in the handout.
\end{laterbox}

\subsection{Q3. $-\vv=(-1)\vv$}

\begin{promptbox}
Prove that $-\vv=(-1)\vv$.
\end{promptbox}

\begin{coursebox}
It is enough to check that $(-1)\vv$ is an additive inverse of $\vv$, because
additive inverses are unique (if $\vv+\ww=\zero=\vv+\vec{u}$ then
$\ww=\zero+\ww=(-\vv+\vv)+\ww=-\vv+(\vv+\ww)=-\vv+\zero=-\vv$, and likewise
$\vec{u}=-\vv$).

Now
\[
\vv+(-1)\vv=1\cdot\vv+(-1)\vv=\bigl(1+(-1)\bigr)\vv=0\vv=\zero,
\]
using Q2. Therefore $-\vv=(-1)\vv$.
\end{coursebox}

\begin{laterbox}
The coordinate proof is again $[\vv]+[(-1)\vv]=[\vv]+(-1)[\vv]=\zero$ in $\F^n$.
The axiom-level proof is the one you want on a midterm at this stage of the
course.
\end{laterbox}

\subsection{Q4. Cancellation / $\vv+\ww=\vv$ implies $\ww=\zero$}

\begin{promptbox}
Show that if $\vv+\ww=\vv$, then $\ww=\zero$.
\end{promptbox}

\begin{coursebox}
Add $-\vv$ to both sides:
\[
-\vv+(\vv+\ww)=-\vv+\vv
\implies
(-\vv+\vv)+\ww=\zero
\implies
\zero+\ww=\zero
\implies
\ww=\zero.
\]
\end{coursebox}

\begin{laterbox}
This is the cancellation law in the abelian group $(V,+)$. Equivalently,
the translation map $T_{\vv}(\ww)=\vv+\ww$ is injective (in fact invertible,
with inverse $T_{-\vv}$).
\end{laterbox}

\subsection{Q5. $\vv=-\vv$ need not imply $\vv=\zero$}

\begin{promptbox}
Show that $\vv=-\vv$ does not necessarily imply $\vv=\zero$.
\end{promptbox}

\begin{coursebox}
The equation $\vv=-\vv$ is equivalent to $2\vv=\zero$. If $1+1\neq 0$ in $\F$
(characteristic not $2$), we may multiply by the scalar $2^{-1}$ and conclude
$\vv=\zero$. So a counterexample must live in characteristic $2$.

Take $\F=\F_2=\{0,1\}$ with $1+1=0$, and take $V=\F_2$ as a vector space over
itself. Let $\vv=1$. Then
\[
-\vv=(-1)\vv=1\cdot 1=1=\vv,
\]
but $\vv\neq\zero$.

(The same phenomenon occurs in any vector space over a field of characteristic
$2$: every vector satisfies $\vv=-\vv$.)
\end{coursebox}

\begin{laterbox}
Once you have bases: if $\operatorname{char}\F=2$ and $\dim V\ge 1$, any nonzero
basis vector is a counterexample. If $\operatorname{char}\F\neq 2$, then
$\vv=-\vv$ really does imply $\vv=\zero$, so the hypothesis ``not necessarily''
is a characteristic-$2$ statement, not a dimension statement.
\end{laterbox}
